% §9 Ethical Considerations — target 0.25 page (mandatory per ACSAC CFP)

\section{Ethical Considerations}
\label{sec:ethical}

\paragraph{Isolation}
All experiments run on an isolated Proxmox cluster with a dedicated bridge
(\texttt{vmbr1}) and no route to external networks; the OpenWrt gateway's
\texttt{WAN} interface is firewall-isolated. No live hosts, production
services, or third-party infrastructure are probed.

\paragraph{Responsible disclosure}
\systemname{} vulnerabilities are deterministically injected into our own
lab VMs via Ansible. No 0-days are introduced; all categories (misconfig,
weak credentials, Terrapin CVE-2023-48795, outdated Dropbear) correspond to
patterns already publicly documented. No responsible-disclosure process is
triggered.

\paragraph{Misuse risk of released artifact}
The benchmark artifact (Ansible playbooks, ground truth, evaluator) does not
embed exploitation payloads for unknown vulnerabilities, and the injected
weaknesses are common misconfigurations documented in OWASP~IoT~Top~10.
The harness uses only standard pentest tools already in widespread use.
We judge the net societal benefit — reproducible evaluation of LLM pentest
agents, and a concrete evaluation surface for future defenders — to exceed
the incremental misuse risk.

\paragraph{IRB / human subjects}
No human subjects are involved. No personal data is collected or processed.

\paragraph{LLM costs and environmental footprint}
All reported experiments together consume an estimated \todo{X} kWh and
cost under \todo{USD Y} in API billing.
